\documentclass[../../main.tex]{subfiles}% 注意这里的文件路径不能用 ./main.tex ,否则用latexmk编译子文件会报错
\graphicspath{{\subfix{./image/}}{\subfix{../../image/}}} % 指定图片引用路径,后续可以直接使用图片文件名
% 如果图片存放在上级目录,则使用 ../../image/ 作为备用路径

% 例如:
% \begin{figure}[H]
% \centering
% \includegraphics[width=0.3\textwidth]{图.png}
% \caption{}
% \label{figure:图}
% \end{figure}
% 注意:上述\label{}一定要放在\caption{}之后,否则引用图片序号会只会显示??.

\begin{document}

\section{整函数与亚纯函数}

\begin{theorem}\label{theorem:定理5.3.1}
设$f$为一整函数,则
\begin{enumerate}[(1)]
\item\label{theorem:定理5.3.1-1} $z = \infty$为$f$的解析点或可去奇点的充要条件是$f$一定是常数.

\item\label{theorem:定理5.3.1-2} $z = \infty$为\( f \) 的一个 \( m \) 阶极点的充要条件是\( f \) 是一个 \( m \) 次多项式.
\end{enumerate}
\end{theorem}
\begin{proof}
\begin{enumerate}[(1)]
\item 充分性显然，只证必要性.
由于 \(f\) 在整个复平面 \(\mathbb{C}\) 上全纯，即 \(f\) 为整函数，则由 \cref{theorem:全纯函数的Taylor展开} 知 \(f\) 在 \(\mathbb{C}\) 上有 Taylor 展开式
\begin{align}
f(z) = \sum_{n = 0}^{\infty} a_n z^n. \label{eq:taylor_whole}
\end{align}
下面根据无穷远点的不同情形分情况讨论：
\begin{enumerate}[(i)]
\item 若无穷远点为 \(f\) 的可去奇点，根据 \subcref{theorem:无穷远点为可去奇点的充要条件}{theorem:无穷远点为可去奇点的充要条件-1}，在无穷远点的去心邻域 \(R<|z|<+\infty\) 内，\(f\) 的 Laurent 展开式只含有常数项和负幂次项，即无正幂次项.
\item 若 \(f\) 在无穷远点处全纯，这等价于函数 \(g(w)=f\left(\frac{1}{w}\right)\) 在 \(w=0\) 处全纯. 因而 \(g(w)\) 在 \(0\) 的某邻域内有 Taylor 展开式 \(g(w)=\sum_{n=0}^{\infty} c_n w^n\). 代入\(z=\frac{1}{w}\) 可得
\begin{align*}
f(z)=g\left( \frac{1}{z} \right) =\sum_{n=0}^{\infty}{c_nz^{-n}},\quad R<|z|<+\infty .
\end{align*}
这说明 \(f\) 在无穷远点去心邻域内的 Laurent 展开同样不含正幂次项.
\end{enumerate}
综合上述两种情形可知，在无穷远点的去心邻域 \(R<|z|<+\infty\) 内，\(f\) 的 Laurent 展开式中所有 \(z\) 的正幂次项系数必须全部为零. 再结合\eqref{eq:taylor_whole}式和Laurent展开式的唯一性可知
\begin{align*}
a_1 = a_2 = \cdots = 0.
\end{align*}
从而 \(f(z) = a_0\) 恒成立，即 \(f\) 必为常数.

\item 充分性显然,只证必要性.
如果无穷远点是整函数 \( f \) 的一个 \( m \) 阶极点,那么根据\rrefthe{theorem:无穷远点为m阶极点的充要条件}{theorem:无穷远点为m阶极点的充要条件-1},它在无穷远点邻域中的Laurent展开式除去一个 \( m \) 次多项式外只有负次幂的项,因此在展开式\eqref{eq:taylor_whole}中必须有
\[
a_{m + 1} = a_{m + 2} = \cdots = 0,
\]
所以 \( f \) 是一个 \( m \) 次多项式.

\end{enumerate}
\end{proof}

\begin{definition}
不是常数和多项式的整函数称为\textbf{超越整函数}.
\end{definition}
\begin{remark}
如 \( \mathrm{e}^z,\sin z,\cos z \) 等,都是超越整函数.
\end{remark}

\begin{proposition}
无穷远点一定是超越整函数的本性奇点.
\end{proposition}
\begin{proof}
反证,设$f$是超越整函数,若无穷远点不是$f$的本性奇点,则当$f$在无穷远点全纯或无穷远点是$f$的可去奇点时,由\refthe{theorem:定理5.3.1}可知$f$为一个常数.当无穷远点是$f$的$m$阶极点时,由\rrefthe{theorem:定理5.3.1}{theorem:定理5.3.1-2}可知$f$是一个$m$次多项式.这与$f$是超越整函数矛盾!
\end{proof}


\begin{definition}
如果 \( f \) 在整个复平面 \( \mathbb{C} \) 上除去极点外没有其他奇点的单值全纯函数,就称 \( f \) 是一个\textbf{亚纯函数}.
\end{definition}

\begin{proposition}
(1)整函数是亚纯函数.

(2)有理函数
\[
f(z) = \frac{P_n(z)}{Q_m(z)}
\]
也是亚纯函数,这里,\( P_n \) 和 \( Q_m \) 是两个既约的多项式.
\end{proposition}
\begin{proof}

\end{proof}

\begin{proposition}
无穷远点或是有理函数的可去奇点,或是有理函数的极点.
\end{proposition}
\begin{proof}
若记
\[
P_n(z) = a_0 + a_1 z + \cdots + a_n z^n, \, a_n \neq 0,
\]
\[
Q_m(z) = b_0 + b_1 z + \cdots + b_m z^m, \, b_m \neq 0,
\]
那么
\[
f(z) = \frac{P_n(z)}{Q_m(z)} = \frac{1}{z^{m - n}} \frac{a_n + a_{n - 1} \frac{1}{z} + \cdots + a_0 \frac{1}{z^n}}{b_m + b_{m - 1} \frac{1}{z} + \cdots + b_0 \frac{1}{z^m}},
\]
所以
\[
\lim_{z \to \infty} f(z) = \begin{cases} 
\dfrac{a_n}{b_m}, & n = m; \\
\infty, & n > m; \\
0, & n < m.
\end{cases}
\]
这说明 \( z = \infty \) 或是 \( f \) 的可去奇点,或是 \( f \) 的极点.
\end{proof}

\begin{theorem}\label{theorem:定理5.3.3}
若 \( z = \infty \) 是亚纯函数 \( f \) 的可去奇点或极点,则 \( f \) 一定是有理函数.
\end{theorem}
\begin{proof}
因 \( z = \infty \) 是 \( f \) 的可去奇点或极点,故必存在 \( R > 0 \),使得 \( f \) 在 \( R < |z| < \infty \) 中全纯.在 \( |z| \leqslant R \) 中,\( f \) 最多只能有有限个极点.因若有无穷多个极点 \( z_j, j = 1, 2, \cdots \),则 \( \{ z_j \} \) 必有收敛的子列 \( \{ z_{k_j} \} \),设其极限为 \( a \),则 \( |a| \leqslant R \),显然若$a$是极点,则\( a \)不是孤立奇点,矛盾!若$a$不是极点,则由$f$是亚纯函数可知$f$在$a$处全纯,但$f$在$a$的任意邻域内必有极点,矛盾!今设 \( z_1, \cdots, z_n \) 为 \( f \) 在 \( |z| \leqslant R \) 中的有限个极点,它们的阶分别为 \( m_1, \cdots, m_n \).由\rrefthe{theorem:m阶极点的充要条件}{theorem:m阶极点的充要条件-1}可知\( f \) 在 \( z_j (j = 1, \cdots, n) \) 附近的Laurent展开的主要部分为
\[
h_j(z) = \frac{c_{-1}^{(j)}}{z - z_j} + \frac{c_{-2}^{(j)}}{(z - z_j)^2} + \cdots + \frac{c_{-m_j}^{(j)}}{(z - z_j)^{m_j}}.
\]
设 \( f \) 在 \( \infty \) 的邻域内的Laurent展开的主要部分为 \( g \),由\rrefthe{theorem:无穷远点为m阶极点的充要条件}{theorem:无穷远点为m阶极点的充要条件-1}和\rrefthe{theorem:无穷远点为可去奇点的充要条件}{theorem:无穷远点为可去奇点的充要条件-1}可知,当 \( z = \infty \) 是 \( f \) 的极点时,\( g \) 是一个多项式;当 \( z = \infty \) 是 \( f \) 的可去奇点时,\( g \equiv 0 \).令
\[
F(z) = f(z) - h_1(z) - \cdots - h_n(z) - g(z),
\]
显然,\( F \) 在 \( \mathbb{C}_{\infty} \) 中除 \( z_1, \cdots, z_n \) 和 \( \infty \) 外是全纯的,而在 \( z_1, \cdots, z_n \) 和 \( \infty \) 这些点上,\( f \) 的主要部分都已经被消去,因而也是全纯的.所以,\( F \) 是 \( \mathbb{C}_{\infty} \) 上的全纯函数,因而由\refthe{theorem:定理5.3.1},\( F \) 是一个常数 \( c \).于是
\[
f(z) = c + g(z) + \sum_{j = 1}^{n} h_j(z),
\]
所以 \( f \) 是有理函数.
\end{proof}
\begin{remark}
这里,我们顺便得到了这样一个结论:\textbf{任何有理函数一定能分解成部分分式之和,而且这种分解是唯一的.}这个结论在计算有理函数的不定积分时已经被多次用过.
\end{remark}

\begin{definition}
非有理函数的亚纯函数称为\textbf{超越亚纯函数}.
\end{definition}

\begin{theorem}\label{theorem:定理5.3.4}
\( \mathrm{Aut}(\mathbb{C}) \) 由所有的一次多项式组成.
\end{theorem}
\begin{proof}
设 \( f(z) = az + b, a \neq 0 \),则显然 \( f \in \mathrm{Aut}(\mathbb{C}) \).反之,对于任意的 \( f \in \mathrm{Aut}(\mathbb{C}) \),因为 \( f \) 是整函数,如果 \( \infty \) 是它的可去奇点,则由\refthe{theorem:定理5.3.1},\( f \) 是一个常数,这不可能.如果 \( \infty \) 是 \( f \) 的本性奇点,则由\refthe{theorem:Weierstrass-定理5.2.5},对于任意 \( A \in \mathbb{C} \),必有 \( z_n \to \infty \),使得 \( \lim\limits_{n \to \infty} f(z_n) = A \).现在记 \( f(z_n) = w_n \),则 \( z_n = f^{-1}(w_n) \),两端令 \( n \to \infty \),即得 \( f^{-1}(A) = \infty \).这说明 \( A \) 是 \( f^{-1} \) 的一个极点,与 \( f^{-1} \) 是整函数相矛盾.由此可知 \( \infty \) 必为 \( f \) 的极点,由\rrefthe{theorem:定理5.3.1}{theorem:定理5.3.1-2}知道,\( f \) 是一个多项式.又因为 \( f \) 在 \( \mathbb{C} \) 上是单叶的,所以 \( f \) 只能是一次多项式.
\end{proof}

\begin{theorem}\label{theorem:定理5.3.5}
\( \mathrm{Aut}(\mathbb{C}_{\infty}) \) 由所有的分式线性变换组成.
\end{theorem}
\begin{proof}
因为是在 \( \mathbb{C}_{\infty} \) 上讨论,\( \mathrm{Aut}(\mathbb{C}_{\infty}) \) 中的元素不再是全纯函数,而是亚纯函数.由\hyperref[section:2.2.5]{分式线性变换}的讨论知道,任何分式线性变换都是 \( \mathrm{Aut}(\mathbb{C}_{\infty}) \) 中的元素.现设 \( f \in \mathrm{Aut}(\mathbb{C}_{\infty}) \),则 \( f \) 必为亚纯函数,而且 \( \infty \) 必是 \( f \) 的可去奇点或极点.由\refthe{theorem:定理5.3.3},\( f \) 必为有理函数,再由它的单叶性,它只能是分式线性变换.
\end{proof}

\begin{example}
\begin{enumerate}[(1)]
\item 设 \(f(z)\) 是亚纯函数, 证明: \(f(z)\) 的任意阶零点和任意阶极点均为 \(\frac{f'(z)}{f(z)}\) 的一阶极点.

\item 设 \(C\) 是一条周线(简单闭曲线), \(f(z)\) 在 \(C\) 的内部区域亚纯, 在 \(C\) 上解析且不为零. 证明:
\begin{align*}
\frac{1}{2\pi \mathrm{i}} \int_C \frac{f'(z)}{f(z)} \mathrm{d}z = N(f, C) - P(f, C).
\end{align*}
其中 \(N(f, C)\) 是 \(f(z)\) 在 \(C\) 内部的零点个数(计重数), \(P(f, C)\) 是极点个数(计重数).
\end{enumerate}
\end{example}
\begin{proof}
\begin{enumerate}[(1)]
\item 设 \(z_0\) 是 \(f(z)\) 的 \(m\) 阶零点 (\(m \geqslant 1\)), 则在 \(z_0\) 的邻域内 \(f(z) = (z-z_0)^m \varphi(z)\), 其中 \(\varphi(z)\) 在 \(z_0\) 处解析且 \(\varphi(z_0) \neq 0\).
求导数得 \(f'(z) = m(z-z_0)^{m-1}\varphi(z) + (z-z_0)^m \varphi'(z)\), 于是
\begin{align*}
\frac{f^{\prime}\left( z \right)}{f\left( z \right)}=\frac{m}{z-z_0}+\frac{\varphi ^{\prime}\left( z \right)}{\varphi \left( z \right)}=\frac{m\varphi \left( z \right) +\left( z-z_0 \right) \varphi ^{\prime}\left( z \right)}{\left( z-z_0 \right) \varphi \left( z \right)}.
\end{align*}
显然$z_0$是$\frac{f(z)}{f'(z)}$的一阶零点,故\(z_0\) 是 \(\frac{f'(z)}{f(z)}\)的一阶极点.又因为 \(\frac{\varphi'(z)}{\varphi(z)}\) 在 \(z_0\) 处解析, 所以有
\begin{align*}
\mathrm{Res}\left( \frac{f^{\prime}\left( z \right)}{f\left( z \right)},z_0 \right) =\lim\limits_{z \to z_0} (z-z_0)\frac{f'(z)}{f(z)} = m \neq 0.
\end{align*}


设 \(z_p\) 是 \(f(z)\) 的 \(k\) 阶极点 (\(k \geqslant 1\)), 则在 \(z_p\) 的去心邻域内 \(f(z) = (z-z_p)^{-k} \psi(z)\), 其中 \(\psi(z)\) 在 \(z_p\) 处解析且 \(\psi(z_p) \neq 0\).

求导数得 \(f'(z) = -k(z-z_p)^{-k-1}\psi(z) + (z-z_p)^{-k} \psi'(z)\), 于是
\begin{align*}
\frac{f^{\prime}\left( z \right)}{f\left( z \right)}=\frac{-k}{z-z_p}+\frac{\psi' \left( z \right)}{\psi \left( z \right)}=\frac{-k\psi \left( z \right) +\left( z-z_p \right) \psi ^{\prime}\left( z \right)}{\left( z-z_p \right) \psi \left( z \right)}.
\end{align*}
显然$z_0$是$\frac{f(z)}{f'(z)}$的一阶零点,故\(z_0\) 是 \(\frac{f'(z)}{f(z)}\)的一阶极点.又因为 \(\frac{\psi'(z)}{\psi(z)}\) 在 \(z_p\) 处解析, 所以有
\begin{align*}
\mathrm{Res}\left( \frac{f^{\prime}\left( z \right)}{f\left( z \right)},z_p \right) =\lim\limits_{z \to z_p} (z-z_p)\frac{f'(z)}{f(z)} = -k \neq 0,
\end{align*}

\item 设 \(z_1, z_2, \cdots, z_s\) 为 \(f(z)\) 在 \(C\) 内部的全部零点, 重数分别为 \(m_1, m_2, \cdots, m_s\); 设 \(p_1, p_2, \cdots, p_t\) 为 \(f(z)\) 在 \(C\) 内部的全部极点, 阶数分别为 \(k_1, k_2, \cdots, k_t\).
由(1)的结论可知, \(z_i\) 为 \(\frac{f'(z)}{f(z)}\) 的一阶极点, 留数为 \(m_i\); \(p_j\) 为 \(\frac{f'(z)}{f(z)}\) 的一阶极点, 留数为 \(-k_j\).
由于 \(f(z)\) 在 \(C\) 上解析且不为零, 所以 \(\frac{f'(z)}{f(z)}\) 在 \(C\) 上解析, 从而 \(\frac{f'(z)}{f(z)}\) 在 \(C\) 内部除上述点外无其他奇点. 根据\hyperref[theorem:留数定理(残数定理)-定理5.4.9]{留数定理}, 有
\begin{align*}
\frac{1}{2\pi \mathrm{i}} \int_C \frac{f'(z)}{f(z)} \mathrm{d}z = \sum\limits_{i=1}^{s} \mathrm{Res}\left(\frac{f'}{f}, z_i\right) + \sum\limits_{j=1}^{t} \mathrm{Res}\left(\frac{f'}{f}, p_j\right),
\end{align*}
即
\begin{align*}
\frac{1}{2\pi \mathrm{i}} \int_C \frac{f'(z)}{f(z)} \mathrm{d}z = \sum\limits_{i=1}^{s} m_i - \sum\limits_{j=1}^{t} k_j.
\end{align*}
根据定义, \(\sum\limits_{i=1}^{s} m_i = N(f, C)\), \(\sum\limits_{j=1}^{t} k_j = P(f, C)\), 故
\begin{align*}
\frac{1}{2\pi \mathrm{i}} \int_C \frac{f'(z)}{f(z)} \mathrm{d}z = N(f, C) - P(f, C).
\end{align*}
\end{enumerate}
\end{proof}







\end{document}